% GENERATED FILE. Edit canonical lessons, then run npm run generate:books.
% canonical-source-hash: fnv1a64:711af0f98f12d6ed
% canonical-lessons: AR-C465-taqa, AR-C465-ghaz, AR-C465-naft, AR-C465-fahm, AR-C465-istahlaka

\chapter{Energy}
\label{ch:ar-energy}
\hlchaptermodality{465}

\begin{chapteropening}
\textbf{By the end of this chapter:} \emph{I can say energy, gas, crude oil, coal and to consume.}

Complete the last lesson of chapter 465: I can say energy, gas, crude oil, coal and to consume.
\end{chapteropening}

\section[\texorpdfstring{\ar{طاقة} (ṭāqa)}{طاقة (ṭāqa)}]{\ar{طاقة} (ṭāqa) — energy; power}

\label{lesson:AR-C465-taqa}



\begin{quote}
Before the new one: say the Arabic for to pollute, then the Arabic for to sort.
\end{quote}

\subsection*{You'll want to know: \ar{طاقة}}
\textbf{\ar{طاقة}} — \emph{ṭāqa} — \textquotedblleft{}energy; power\textquotedblright{}.

\begin{quote}
\ar{نستعمل} \ar{طاقة} \ar{الشمس} \ar{في} \ar{منزلنا}.
\end{quote}

\begin{quote}
\emph{nastaʿmilu ṭāqat ash-shams fī manzilinā.}
\end{quote}

\begin{quote}
We use the sun's energy in our house.
\end{quote}

\begin{cousinweb}
\textbf{\ar{طاقة}} (\emph{ṭāqa}) is built on the native Arabic root \emph{ṭ-w-q}, \textquotedblleft{}to be able\textquotedblright{} (Lane). A \emph{ṭāqa} was first strength, what one is able to do.
\end{cousinweb}

\begin{grammarlens}[title={What you've built}]
One more word for getting around.
\end{grammarlens}

\subsection*{Your turn}
\begin{itemize}
\raggedright
  \item \emph{Say it:} \emph{ṭāqa}
  \item \emph{Say it:} \emph{ṭāqa}, once more
  \item \emph{Say it:} \emph{nastaʿmilu ṭāqat ash-shams fī manzilinā.} — We use the sun's energy in our house.
  \item \emph{Recall:} say the Arabic for to pollute, then the Arabic for to sort, then say \emph{ṭāqa} again
\end{itemize}

\begin{tcolorbox}[breakable,colback=teal!4,colframe=teal!35!black,title={Before you move on}]
What does \ar{طاقة} mean? (\textquotedblleft{}Energy; power\textquotedblright{}.) Say it once more.
\end{tcolorbox}
\section[\texorpdfstring{\ar{غاز} (ghāz)}{غاز (ghāz)}]{\ar{غاز} (ghāz) — gas}

\label{lesson:AR-C465-ghaz}



\begin{quote}
Before the new one: say the Arabic for to sort, then the Arabic for energy.
\end{quote}

\subsection*{You'll want to know: \ar{غاز}}
\textbf{\ar{غاز}} — \emph{ghāz} — \textquotedblleft{}gas\textquotedblright{}.

\begin{quote}
\ar{نطبخ} \ar{على} \ar{الغاز}.
\end{quote}

\begin{quote}
\emph{naṭbukhu ʿalā al-ghāz.}
\end{quote}

\begin{quote}
We cook on gas.
\end{quote}

\begin{cousinweb}
\textbf{\ar{غاز}} (\emph{ghāz}) was borrowed from English \emph{gas} (EtymDB).
\end{cousinweb}

\begin{grammarlens}[title={What you've built}]
One more word for getting around.
\end{grammarlens}

\subsection*{Your turn}
\begin{itemize}
\raggedright
  \item \emph{Say it:} \emph{ghāz}
  \item \emph{Say it:} \emph{ghāz}, once more
  \item \emph{Say it:} \emph{naṭbukhu ʿalā al-ghāz.} — We cook on gas.
  \item \emph{Recall:} say the Arabic for to sort, then the Arabic for energy, then say \emph{ghāz} again
\end{itemize}

\begin{tcolorbox}[breakable,colback=teal!4,colframe=teal!35!black,title={Before you move on}]
What does \ar{غاز} mean? (\textquotedblleft{}Gas\textquotedblright{}.) Say it once more.
\end{tcolorbox}
\section[\texorpdfstring{\ar{نفط} (nafṭ)}{نفط (nafṭ)}]{\ar{نفط} (nafṭ) — crude oil; petroleum}

\label{lesson:AR-C465-naft}



\begin{quote}
Before the new one: say the Arabic for energy, then the Arabic for gas.
\end{quote}

\subsection*{You'll want to know: \ar{نفط}}
\textbf{\ar{نفط}} — \emph{nafṭ} — \textquotedblleft{}crude oil; petroleum\textquotedblright{}.

\begin{quote}
\ar{في} \ar{هذه} \ar{المنطقة} \ar{آبار} \ar{نفط}.
\end{quote}

\begin{quote}
\emph{fī hādhihi al-minṭaqa ābār nafṭ.}
\end{quote}

\begin{quote}
There are oil wells in this area.
\end{quote}

\begin{cousinweb}
\textbf{\ar{نفط}} (\emph{nafṭ}) is a very old word for oil; the origin is not settled.
\end{cousinweb}

\begin{grammarlens}[title={What you've built}]
One more word for getting around.
\end{grammarlens}

\subsection*{Your turn}
\begin{itemize}
\raggedright
  \item \emph{Say it:} \emph{nafṭ}
  \item \emph{Say it:} \emph{nafṭ}, once more
  \item \emph{Say it:} \emph{fī hādhihi al-minṭaqa ābār nafṭ.} — There are oil wells in this area.
  \item \emph{Recall:} say the Arabic for energy, then the Arabic for gas, then say \emph{nafṭ} again
\end{itemize}

\begin{tcolorbox}[breakable,colback=teal!4,colframe=teal!35!black,title={Before you move on}]
What does \ar{نفط} mean? (\textquotedblleft{}Crude oil; petroleum\textquotedblright{}.) Say it once more.
\end{tcolorbox}
\section[\texorpdfstring{\ar{فحم} (faḥm)}{فحم (faḥm)}]{\ar{فحم} (faḥm) — coal; charcoal}

\label{lesson:AR-C465-fahm}



\begin{quote}
Before the new one: say the Arabic for gas, then the Arabic for crude oil.
\end{quote}

\subsection*{You'll want to know: \ar{فحم}}
\textbf{\ar{فحم}} — \emph{faḥm} — \textquotedblleft{}coal; charcoal\textquotedblright{}.

\begin{quote}
\ar{شوى} \ar{أبي} \ar{اللحم} \ar{على} \ar{الفحم}.
\end{quote}

\begin{quote}
\emph{shawā abī al-laḥm ʿalā al-faḥm.}
\end{quote}

\begin{quote}
My father grilled the meat over charcoal.
\end{quote}

\begin{cousinweb}
\textbf{\ar{فحم}} (\emph{faḥm}) is built on the native Arabic root \emph{f-ḥ-m}, \textquotedblleft{}charcoal\textquotedblright{} (Lane). EtymDB lists a related Hebrew word for coal.
\end{cousinweb}

\begin{grammarlens}[title={What you've built}]
One more word for getting around.
\end{grammarlens}

\subsection*{Your turn}
\begin{itemize}
\raggedright
  \item \emph{Say it:} \emph{faḥm}
  \item \emph{Say it:} \emph{faḥm}, once more
  \item \emph{Say it:} \emph{shawā abī al-laḥm ʿalā al-faḥm.} — My father grilled the meat over charcoal.
  \item \emph{Recall:} say the Arabic for gas, then the Arabic for crude oil, then say \emph{faḥm} again
\end{itemize}

\begin{tcolorbox}[breakable,colback=teal!4,colframe=teal!35!black,title={Before you move on}]
What does \ar{فحم} mean? (\textquotedblleft{}Coal; charcoal\textquotedblright{}.) Say it once more.
\end{tcolorbox}
\section[\texorpdfstring{\ar{استهلك} (istahlaka)}{استهلك (istahlaka)}]{\ar{استهلك} (istahlaka) — to consume; to use up}

\label{lesson:AR-C465-istahlaka}



\begin{quote}
Before the new one: say the Arabic for crude oil, then the Arabic for coal.
\end{quote}

\subsection*{You'll want to know: \ar{استهلك}}
\textbf{\ar{استهلك}} — \emph{istahlaka} — \textquotedblleft{}to consume; to use up\textquotedblright{}.

\begin{quote}
\ar{تستهلك} \ar{السيارة} \ar{الكبيرة} \ar{بنزينا} \ar{كثيرا}.
\end{quote}

\begin{quote}
\emph{tastahliku as-sayyāra al-kabīra banzīnan kathīran.}
\end{quote}

\begin{quote}
The big car uses a lot of petrol.
\end{quote}

Present: \ar{يستهلك}, \emph{yastahlik}.

\begin{cousinweb}
\textbf{\ar{استهلك}} (\emph{istahlaka}) is built on the native Arabic root \emph{h-l-k}, \textquotedblleft{}to perish\textquotedblright{} (Lane). It is Form X, \textquotedblleft{}to make perish, to use up\textquotedblright{}.
\end{cousinweb}

\begin{grammarlens}[title={What you've built}]
One more word for getting around.
\end{grammarlens}

\subsection*{Your turn}
\begin{itemize}
\raggedright
  \item \emph{Say it:} \emph{istahlaka}
  \item \emph{Say it:} \emph{istahlaka}, once more
  \item \emph{Say it:} \emph{tastahliku as-sayyāra al-kabīra banzīnan kathīran.} — The big car uses a lot of petrol.
  \item \emph{Recall:} say the Arabic for crude oil, then the Arabic for coal, then say \emph{istahlaka} again
\end{itemize}

\begin{tcolorbox}[breakable,colback=teal!4,colframe=teal!35!black,title={Before you move on}]
What does \ar{استهلك} mean? (\textquotedblleft{}To consume; to use up\textquotedblright{}.) Say it once more.
\end{tcolorbox}
